%%%%%%%%%%%%%%%%%%%%%%%%%%%%%%%%%%%%%%%%%%%%%%%%%%%%%%%%%%%%%%%%%%%%%%%%%%%%%%%
%%  Chapter 10 — CIIR as an External Informational Layer on Quantum Systems
%%
%%  CURRENT THEORY STATE
%%
%%  Definitions:     D3.1–D3.21 (Ch.~3); D5.1–D5.14 (Ch.~5);
%%                   D6.1–D6.8  (Ch.~6); D7.1–D7.9  (Ch.~7);
%%                   D8.1–D8.10 (Ch.~8)
%%  Axioms:          Ax.4.1–Ax.4.6 (Ch.~4)
%%  Proven Theorems: Th.5.1–5.12 (Ch.~5); Th.6.1–6.7 (Ch.~6);
%%                   Th.7.1–7.7 (Ch.~7); Th.8.1–8.9 (Ch.~8);
%%                   Th.9.1–9.6 (Ch.~9)
%%
%%  Dependencies: Chapters 3–9 (Primitive Definitions D3.1–D3.21,
%%                Axioms Ax.4.1–Ax.4.6, Lindblad dynamics Ch.~7,
%%                Measurement theory Ch.~8, Consolidation Ch.~9)
%%
%%  PURPOSE: Embed CIIR as a fixed external informational structure onto
%%           standard quantum-theoretic formalism.  CIIR is NOT modified
%%           or redefined; quantum mechanics is NOT altered.  Instead we
%%           define a rigorous coupling map and derive perturbative
%%           consequences.
%%%%%%%%%%%%%%%%%%%%%%%%%%%%%%%%%%%%%%%%%%%%%%%%%%%%%%%%%%%%%%%%%%%%%%%%%%%%%%%

\chapter{CIIR as an External Informational Layer on Quantum Systems}
\label{ch:quantum_embedding}

\vspace{0.5em}
\noindent
This chapter establishes the formal relationship between the CIIR
framework developed in Chapters~3--9 and standard quantum theory.  The
central design principle is that \emph{neither} structure is modified:
quantum mechanics retains its full Hilbert-space, CPTP, and
measurement axiomatics, while CIIR retains the definitions, axioms,
and theorems proven in the preceding chapters.  We construct an
explicit \emph{coupling map} from the CIIR informational structure
into the algebra of bounded operators on the quantum Hilbert space,
derive perturbative corrections to dynamics and measurement, and
identify rigorous no-go conditions and possible experimental
signatures.

%======================================================================
\section{Quantum Theory as Fixed Base Structure}
\label{sec:ch10:qm_base}
%======================================================================

We begin by stating the axioms of quantum mechanics that serve as the
\emph{unmodified} base layer throughout this chapter.

\begin{axiom}[Hilbert-Space Postulate]
\label{ax:10.1}
A quantum system is associated with a separable complex Hilbert space
$\mathcal{H}$ equipped with inner product
$\langle\cdot|\cdot\rangle$.  Pure states are unit rays in
$\mathcal{H}$; mixed states are density operators
$\rho\in\mathcal{D}(\mathcal{H})$, where
\[
  \mathcal{D}(\mathcal{H})
  \;=\;
  \bigl\{\,\rho\in\CB(\mathcal{H})
         \;\big|\;
         \rho\geq 0,\;\Tr(\rho)=1
  \bigr\}.
\]
\end{axiom}

\begin{axiom}[Born Rule]
\label{ax:10.2}
For a POVM $\{E_k\}_{k\in\mathcal{K}}$ on $\mathcal{H}$ satisfying
$E_k\geq 0$ and $\sum_k E_k = I$, the probability of outcome $k$
given state $\rho$ is
\[
  p_k \;=\; \Tr(E_k\,\rho).
\]
\end{axiom}

\begin{axiom}[Unitary Evolution]
\label{ax:10.3}
Closed-system evolution is governed by a one-parameter family of
unitary operators $\{U(t)\}_{t\in\mathbb{R}}$ generated by a
self-adjoint Hamiltonian $\hatH$:
\[
  \rho(t) \;=\; U(t)\,\rho(0)\,U(t)^{\dagger},
  \qquad
  U(t) = e^{-i\hatH t/\hbar}.
\]
\end{axiom}

\begin{axiom}[Open-System Dynamics — CPTP Maps]
\label{ax:10.4}
The most general physically admissible evolution of a quantum state is
a completely positive, trace-preserving (CPTP) map
$\mathcal{E}:\mathcal{D}(\mathcal{H})\to\mathcal{D}(\mathcal{H})$
admitting a Kraus representation
\[
  \mathcal{E}(\rho)
  \;=\;
  \sum_{i} K_i\,\rho\,K_i^{\dagger},
  \qquad
  \sum_{i} K_i^{\dagger}\,K_i = I.
\]
\end{axiom}

\begin{definition}[Density-Operator Space]
\label{def:10.1}
Let $\mathcal{D}(\mathcal{H})$ denote the convex set of all density
operators on $\mathcal{H}$.  Equipped with the trace norm
$\|\cdot\|_1$, the ambient space
$\mathcal{T}_1(\mathcal{H})$ of trace-class operators is a Banach
space, and $\mathcal{D}(\mathcal{H})\subset\mathcal{T}_1(\mathcal{H})$
is a norm-closed convex subset.
\end{definition}

\begin{definition}[Quantum Channel]
\label{def:10.2}
A \emph{quantum channel} is a linear map
$\mathcal{E}:\mathcal{T}_1(\mathcal{H})\to\mathcal{T}_1(\mathcal{H})$
that is completely positive and trace-preserving.  Equivalently,
$\mathcal{E}$ maps $\mathcal{D}(\mathcal{H})$ into itself.
\end{definition}

\begin{definition}[GKSL Generator]
\label{def:10.3}
The Gorini--Kossakowski--Sudarshan--Lindblad (GKSL) master equation
for a Markovian open quantum system is
\begin{equation}
\label{eq:10.gksl}
  \frac{d\rho}{dt}
  \;=\;
  -\frac{i}{\hbar}\bigl[\hatH,\,\rho\bigr]
  \;+\;
  \sum_{k}
  \Bigl(
    L_k\,\rho\,L_k^{\dagger}
    -
    \tfrac{1}{2}\bigl\{L_k^{\dagger}L_k,\,\rho\bigr\}
  \Bigr),
\end{equation}
where $\hatH=\hatH^{\dagger}$ is the system Hamiltonian and
$\{L_k\}\subset\CB(\mathcal{H})$ are Lindblad operators encoding the
system--environment interaction.  The generator
$\CL_{\mathrm{QM}}(\cdot)$ on the right-hand side is the most general
generator of a quantum-dynamical semigroup on
$\mathcal{D}(\mathcal{H})$.
\end{definition}

\begin{remark}
\label{rem:10.1}
Everything in this section is standard and will \emph{not} be modified
in the sequel.  The CIIR framework enters purely as an additional
external informational constraint, compatible with every axiom stated
above.
\end{remark}

%======================================================================
\section{CIIR as External Informational Constraint System}
\label{sec:ch10:ciir_external}
%======================================================================

We now package the CIIR definitions from Chapters~3--4 into a single
algebraic datum and show that it acts as an external constraint on the
quantum state space without altering the latter's axiomatics.

\begin{definition}[CIIR Informational Structure]
\label{def:10.4}
The \emph{CIIR informational structure} is the tuple
\[
  \mathfrak{I}
  \;=\;
  \bigl(\CM,\,g,\,\kappa,\,\Phi,\,\CA(\HC)\bigr),
\]
where
\begin{enumerate}
  \item $(\CM,g)$ is the constraint manifold with Riemannian metric~$g$
        (Definition~\ref{def:3.7});
  \item $\kappa:\CM\to\mathbb{R}_{\geq 0}$ is the scalar curvature
        function derived from~$g$ (Definition~\ref{def:3.7});
  \item $\Phi:\CR\to\CB(\HC)$ is the interface map
        (Definition~\ref{def:3.18}), mapping the reality space~$\CR$
        (Definition~\ref{def:3.1}) to bounded operators on the
        cognitive state space~$\HC$ (Definition~\ref{def:3.11});
  \item $\CA(\HC)\subseteq\CB(\HC)$ is the observable algebra
        (Definition~\ref{def:3.20}).
\end{enumerate}
\end{definition}

\begin{definition}[Coupling Map]
\label{def:10.5}
The \emph{coupling map} is the linear map
\[
  C_{\mathfrak{I}}
  :\;
  \Gamma(\CM)
  \;\longrightarrow\;
  \CB(\mathcal{H}),
\]
defined for each smooth section $\sigma\in\Gamma(\CM)$ by
\begin{equation}
\label{eq:10.coupling}
  C_{\mathfrak{I}}(\sigma)
  \;=\;
  \int_{\CM}
  \kappa(m)\;\Phi\bigl(\iota(m)\bigr)\;
  \sigma(m)\;
  d\mu_g(m),
\end{equation}
where $\iota:\CM\hookrightarrow\CR$ is the canonical embedding and
$d\mu_g$ is the Riemannian volume form on $(\CM,g)$.
\end{definition}

\begin{theorem}[Coupling-Map Boundedness]
\label{thm:10.1}
Let $\kappa_{\max}=\sup_{m\in\CM}|\kappa(m)|<\infty$ and assume that
$\|\Phi\|:=\sup_{r\in\CR}\|\Phi(r)\|_{\mathrm{op}}<\infty$
(guaranteed by Axiom~\ref{ax:4.3}).  Then for every
$\sigma\in\Gamma(\CM)$ with $\|\sigma\|_{L^1(\CM,\mu_g)}<\infty$,
\[
  \bigl\|C_{\mathfrak{I}}(\sigma)\bigr\|_{\mathrm{op}}
  \;\leq\;
  \kappa_{\max}\cdot\|\Phi\|\cdot\|\sigma\|_{L^1(\CM,\mu_g)}.
\]
In particular, $C_{\mathfrak{I}}(\sigma)\in\CB(\mathcal{H})$.
\end{theorem}

\begin{proof}
By the triangle inequality for the operator norm and the sub-multiplicativity
of the integrand,
\begin{align*}
  \bigl\|C_{\mathfrak{I}}(\sigma)\bigr\|_{\mathrm{op}}
  &\;\leq\;
  \int_{\CM}
  |\kappa(m)|\;\bigl\|\Phi\bigl(\iota(m)\bigr)\bigr\|_{\mathrm{op}}\;
  |\sigma(m)|\;d\mu_g(m) \\
  &\;\leq\;
  \kappa_{\max}\cdot\|\Phi\|
  \int_{\CM}|\sigma(m)|\;d\mu_g(m) \\
  &\;=\;
  \kappa_{\max}\cdot\|\Phi\|\cdot\|\sigma\|_{L^1(\CM,\mu_g)}.
\end{align*}
Since $\kappa_{\max}<\infty$ and $\|\Phi\|<\infty$ by hypothesis, the
result follows.  Boundedness of $\|\Phi\|$ is ensured by
Axiom~\ref{ax:4.3} (Interface Completeness).
\end{proof}

\begin{corollary}[State-Space Preservation]
\label{cor:10.1}
The image of $C_{\mathfrak{I}}$ lies in $\CB(\mathcal{H})$.  In
particular, $C_{\mathfrak{I}}$ does not map out of the operator algebra
on the quantum Hilbert space; hence the quantum state space
$\mathcal{D}(\mathcal{H})$ and its convex structure are preserved
under CIIR coupling.
\end{corollary}

\begin{proof}
Immediate from Theorem~\ref{thm:10.1}: every element in the image is
a bounded operator, and no modification of $\mathcal{D}(\mathcal{H})$
is induced since $C_{\mathfrak{I}}$ maps into $\CB(\mathcal{H})$, not
into $\mathcal{D}(\mathcal{H})$ itself.
\end{proof}

\begin{figure}[t]
\centering
\fbox{\parbox{0.85\linewidth}{%
\textbf{Figure~10.1: Coupling Architecture}\\[4pt]
Schematic of the coupling map $C_{\mathfrak{I}}$.  The CIIR
informational structure $\mathfrak{I}=(\CM,g,\kappa,\Phi,\CA(\HC))$
(left) maps via $C_{\mathfrak{I}}$ into the algebra $\CB(\mathcal{H})$
of bounded operators on the quantum Hilbert space (right).  Quantum
state space $\mathcal{D}(\mathcal{H})\subset\CB(\mathcal{H})$ remains
unmodified.  Arrows indicate: (i)~curvature weighting by~$\kappa$,
(ii)~interface embedding via~$\Phi$, (iii)~integration over~$\CM$.
}}
\caption{Coupling map from CIIR informational structure to quantum
  operator algebra.}
\label{fig:c10:coupling}
\end{figure}

%======================================================================
\section{Coupling Map Between CIIR and Quantum Channels}
\label{sec:ch10:channel_coupling}
%======================================================================

We now construct the CIIR-modulated quantum channel and prove that it
preserves complete positivity and trace preservation.

\begin{definition}[Standard CPTP Channel]
\label{def:10.6}
A standard CPTP channel $\mathcal{E}_{\mathrm{QM}}$ acts on
$\rho\in\mathcal{D}(\mathcal{H})$ via Kraus operators
$\{K_i\}_{i\in\mathcal{I}}$:
\begin{equation}
\label{eq:10.cptp}
  \mathcal{E}_{\mathrm{QM}}(\rho)
  \;=\;
  \sum_{i\in\mathcal{I}} K_i\,\rho\,K_i^{\dagger},
  \qquad
  \sum_{i\in\mathcal{I}} K_i^{\dagger}\,K_i = I.
\end{equation}
\end{definition}

\begin{definition}[Constraint Lindblad Operators]
\label{def:10.7}
For each index $k$ in a finite set $\mathcal{K}_{\mathfrak{I}}$,
the \emph{constraint Lindblad operator} derived from the CIIR
informational structure is
\begin{equation}
\label{eq:10.constraint_lindblad}
  L_k^{(\mathfrak{I})}
  \;=\;
  \sqrt{\kappa(m_k)}\;
  \Phi\bigl(\iota(m_k)\bigr),
\end{equation}
where $\{m_k\}_{k\in\mathcal{K}_{\mathfrak{I}}}\subset\CM$ is a
finite sampling of the constraint manifold at points where the
curvature $\kappa$ is non-vanishing, and $\Phi$ is the interface map
(Definition~\ref{def:3.18}).
\end{definition}

\begin{definition}[CIIR Dissipative Correction]
\label{def:10.8}
The \emph{CIIR dissipative correction} is the superoperator
\begin{equation}
\label{eq:10.delta_I}
  \Delta_{\mathfrak{I}}(\rho)
  \;=\;
  \varepsilon
  \sum_{k\in\mathcal{K}_{\mathfrak{I}}}
  \Bigl(
    L_k^{(\mathfrak{I})}\,\rho\,L_k^{(\mathfrak{I})\dagger}
    -
    \tfrac{1}{2}
    \bigl\{
      L_k^{(\mathfrak{I})\dagger}L_k^{(\mathfrak{I})},\;\rho
    \bigr\}
  \Bigr),
\end{equation}
where $\varepsilon>0$ is a dimensionless coupling parameter
(specified in Section~\ref{sec:ch10:perturbative}).
\end{definition}

\begin{definition}[CIIR-Modulated Channel]
\label{def:10.9}
The \emph{CIIR-modulated channel} is
\begin{equation}
\label{eq:10.ciir_channel}
  \mathcal{E}_{\mathrm{CIIR}}(\rho)
  \;=\;
  \mathcal{E}_{\mathrm{QM}}(\rho)
  \;+\;
  \Delta_{\mathfrak{I}}(\rho).
\end{equation}
\end{definition}

\begin{theorem}[CPTP Preservation]
\label{thm:10.2}
Suppose $\mathcal{E}_{\mathrm{QM}}$ is a CPTP channel
(Definition~\ref{def:10.6}) and $\Delta_{\mathfrak{I}}$ is the CIIR
dissipative correction (Definition~\ref{def:10.8}).  Then
$\mathcal{E}_{\mathrm{CIIR}}$ as defined in
equation~\eqref{eq:10.ciir_channel} is a valid CPTP map on
$\mathcal{D}(\mathcal{H})$.
\end{theorem}

\begin{proof}
\emph{Trace preservation.}  Since $\Delta_{\mathfrak{I}}$ is in
Lindblad form, it is trace-annihilating:
\[
  \Tr\bigl(\Delta_{\mathfrak{I}}(\rho)\bigr)
  \;=\;
  \varepsilon\sum_k
  \Bigl(
    \Tr\bigl(L_k^{(\mathfrak{I})}\rho\,L_k^{(\mathfrak{I})\dagger}\bigr)
    -
    \Tr\bigl(L_k^{(\mathfrak{I})\dagger}L_k^{(\mathfrak{I})}\rho\bigr)
  \Bigr)
  \;=\;0,
\]
by cyclicity of the trace.  Hence
$\Tr\bigl(\mathcal{E}_{\mathrm{CIIR}}(\rho)\bigr)
=\Tr\bigl(\mathcal{E}_{\mathrm{QM}}(\rho)\bigr)=1$.

\emph{Complete positivity.}  The map $\mathcal{E}_{\mathrm{CIIR}}$
is the sum of $\mathcal{E}_{\mathrm{QM}}$ and the Lindblad
superoperator $\Delta_{\mathfrak{I}}$.  Rewrite
$\mathcal{E}_{\mathrm{CIIR}}$ as the time-$\delta t$ map of the
GKSL generator
\[
  \CL_{\mathrm{CIIR}}(\rho)
  \;=\;
  \CL_{\mathrm{QM}}(\rho)
  \;+\;
  \varepsilon\sum_k
  \Bigl(
    L_k^{(\mathfrak{I})}\rho\,L_k^{(\mathfrak{I})\dagger}
    -\tfrac{1}{2}\bigl\{L_k^{(\mathfrak{I})\dagger}L_k^{(\mathfrak{I})},\rho\bigr\}
  \Bigr),
\]
which is the sum of two Lindblad generators and hence itself a valid
Lindblad generator~\cite{lindblad1976}.  By the GKSL theorem, the
generated semigroup $\{e^{t\,\CL_{\mathrm{CIIR}}}\}_{t\geq 0}$
consists entirely of CPTP maps.  Setting $\delta t=1$ (in appropriate
units) recovers $\mathcal{E}_{\mathrm{CIIR}}$.
\end{proof}

\begin{lemma}[Correction Norm Bound]
\label{lem:10.1}
For all $\rho\in\mathcal{D}(\mathcal{H})$,
\begin{equation}
\label{eq:10.norm_bound}
  \bigl\|\Delta_{\mathfrak{I}}(\rho)\bigr\|_1
  \;\leq\;
  \varepsilon\cdot\kappa_{\max}\cdot\|\Phi\|
  \cdot
  2\,|\mathcal{K}_{\mathfrak{I}}|,
\end{equation}
where $|\mathcal{K}_{\mathfrak{I}}|$ is the cardinality of the
constraint sampling set.  In particular, if we define
$\|\Phi\|_{\kappa}:= \kappa_{\max}\cdot\|\Phi\|$, then
$\|\Delta_{\mathfrak{I}}(\rho)\|_1 \leq
2\,\varepsilon\,|\mathcal{K}_{\mathfrak{I}}|\,\|\Phi\|_{\kappa}$.
\end{lemma}

\begin{proof}
Each Lindblad term satisfies
$\|L_k^{(\mathfrak{I})}\rho\,L_k^{(\mathfrak{I})\dagger}\|_1
\leq \|L_k^{(\mathfrak{I})}\|_{\mathrm{op}}^2\|\rho\|_1
= \kappa(m_k)\|\Phi(\iota(m_k))\|_{\mathrm{op}}^2$,
with $\|\rho\|_1=1$.  Similarly for the anticommutator term.  Summing
over $k$ and bounding each factor by its supremum yields the result.
\end{proof}

\begin{figure}[t]
\centering
\fbox{\parbox{0.85\linewidth}{%
\textbf{Figure~10.2: Channel Decomposition}\\[4pt]
Diagram showing the CIIR-modulated channel
$\mathcal{E}_{\mathrm{CIIR}} = \mathcal{E}_{\mathrm{QM}} +
\Delta_{\mathfrak{I}}$.  The standard quantum channel
$\mathcal{E}_{\mathrm{QM}}$ (top path) and the CIIR correction
$\Delta_{\mathfrak{I}}$ (bottom path, weighted by~$\varepsilon$)
act in parallel on $\rho$ and their outputs are summed.
The correction magnitude is bounded by
$\varepsilon\cdot\kappa_{\max}\cdot\|\Phi\|$.
}}
\caption{Decomposition of the CIIR-modulated quantum channel.}
\label{fig:c10:channel}
\end{figure}

%======================================================================
\section{Perturbative Dynamics and Consistency Conditions}
\label{sec:ch10:perturbative}
%======================================================================

We now fix the perturbation parameter and derive the leading-order
correction to quantum dynamics imposed by CIIR coupling.

\begin{definition}[Perturbation Parameter]
\label{def:10.10}
The \emph{CIIR perturbation parameter} is the dimensionless ratio
\begin{equation}
\label{eq:10.epsilon}
  \varepsilon
  \;=\;
  \frac{\kappa_{\max}}{E_{\mathrm{typical}}},
\end{equation}
where $E_{\mathrm{typical}}$ is a characteristic energy scale of the
quantum system (e.g., the spectral gap of $\hatH$) and $\kappa_{\max}$
carries units of inverse length squared (curvature).  We assume
$\varepsilon\ll 1$ throughout.
\end{definition}

\begin{theorem}[Perturbative Expansion of State Evolution]
\label{thm:10.3}
Let $\rho(t)$ satisfy the CIIR-modified master equation
\begin{equation}
\label{eq:10.ciir_master}
  \frac{d\rho}{dt}
  \;=\;
  \CL_{\mathrm{QM}}(\rho)
  \;+\;
  \varepsilon\,\CL_{\mathfrak{I}}(\rho),
\end{equation}
where $\CL_{\mathrm{QM}}$ is the standard GKSL generator
(Definition~\ref{def:10.3}) and
\[
  \CL_{\mathfrak{I}}(\rho)
  \;=\;
  \sum_{k\in\mathcal{K}_{\mathfrak{I}}}
  \Bigl(
    L_k^{(\mathfrak{I})}\rho\,L_k^{(\mathfrak{I})\dagger}
    -\tfrac{1}{2}\bigl\{L_k^{(\mathfrak{I})\dagger}L_k^{(\mathfrak{I})},\rho\bigr\}
  \Bigr).
\]
Then $\rho(t)$ admits the asymptotic expansion
\begin{equation}
\label{eq:10.perturbative_expansion}
  \rho(t)
  \;=\;
  \rho_{\mathrm{QM}}(t)
  \;+\;
  \varepsilon\,\rho^{(1)}(t)
  \;+\;
  O(\varepsilon^2),
\end{equation}
where $\rho_{\mathrm{QM}}(t)=e^{t\CL_{\mathrm{QM}}}(\rho_0)$ is the
unperturbed solution and $\rho^{(1)}(t)$ satisfies the inhomogeneous
equation
\begin{equation}
\label{eq:10.first_order}
  \frac{d\rho^{(1)}}{dt}
  \;=\;
  \CL_{\mathrm{QM}}\bigl(\rho^{(1)}\bigr)
  \;+\;
  \CL_{\mathfrak{I}}\bigl(\rho_{\mathrm{QM}}(t)\bigr),
  \qquad
  \rho^{(1)}(0) = 0.
\end{equation}
\end{theorem}

\begin{proof}
Substitute the ansatz~\eqref{eq:10.perturbative_expansion} into
equation~\eqref{eq:10.ciir_master} and collect powers of
$\varepsilon$.

\emph{Order~$\varepsilon^0$:}
$\dot\rho_{\mathrm{QM}} = \CL_{\mathrm{QM}}(\rho_{\mathrm{QM}})$,
which is the standard master equation; satisfied by hypothesis.

\emph{Order~$\varepsilon^1$:}
$\dot\rho^{(1)} = \CL_{\mathrm{QM}}(\rho^{(1)})
+ \CL_{\mathfrak{I}}(\rho_{\mathrm{QM}})$,
which is~\eqref{eq:10.first_order} with initial condition
$\rho^{(1)}(0)=0$ (unperturbed initial state).

The formal solution is given by the Duhamel integral
\[
  \rho^{(1)}(t)
  \;=\;
  \int_0^t
  e^{(t-s)\CL_{\mathrm{QM}}}\,
  \CL_{\mathfrak{I}}\bigl(\rho_{\mathrm{QM}}(s)\bigr)\,ds.
\]
Since $\CL_{\mathrm{QM}}$ generates a contraction semigroup on
$\mathcal{T}_1(\mathcal{H})$ and $\CL_{\mathfrak{I}}$ is bounded
(Lemma~\ref{lem:10.1}), the integrand is norm-continuous and the
integral converges absolutely.  The remainder
$O(\varepsilon^2)$ is controlled by iteration of the Duhamel formula.
\end{proof}

\begin{theorem}[Convergence Condition]
\label{thm:10.4}
The perturbative expansion~\eqref{eq:10.perturbative_expansion}
converges in the trace norm uniformly on compact time intervals
$[0,T]$ provided
\begin{equation}
\label{eq:10.convergence}
  \varepsilon
  \;<\;
  \frac{1}{\bigl\|\CL_{\mathfrak{I}}\bigr\|_{\mathrm{op}}},
  \qquad\text{where}\quad
  \bigl\|\CL_{\mathfrak{I}}\bigr\|_{\mathrm{op}}
  \;\leq\;
  2\sum_{k}\bigl\|L_k^{(\mathfrak{I})}\bigr\|_{\mathrm{op}}^2.
\end{equation}
Moreover, the CIIR coupling $\varepsilon\CL_{\mathfrak{I}}$ is a
bounded perturbation of the generator $\CL_{\mathrm{QM}}$ in the
sense of Kato, so that the perturbed semigroup
$\{e^{t(\CL_{\mathrm{QM}}+\varepsilon\CL_{\mathfrak{I}})}\}_{t\geq 0}$
is a strongly continuous contraction semigroup on
$\mathcal{T}_1(\mathcal{H})$.
\end{theorem}

\begin{proof}
The Lindblad superoperator $\CL_{\mathfrak{I}}$ is bounded with
\[
  \bigl\|\CL_{\mathfrak{I}}(\rho)\bigr\|_1
  \;\leq\;
  2\sum_k\bigl\|L_k^{(\mathfrak{I})}\bigr\|_{\mathrm{op}}^2\;\|\rho\|_1
\]
(triangle inequality plus sub-multiplicativity).  By the Kato
perturbation theorem for strongly continuous semigroups
\cite{kato1995}, if $B$ is a bounded operator and $A$ generates a
$C_0$-semigroup, then $A+B$ also generates a $C_0$-semigroup.  Here
$A=\CL_{\mathrm{QM}}$ (unbounded generator in general) and
$B=\varepsilon\CL_{\mathfrak{I}}$ (bounded).  The perturbation series
converges when $\varepsilon\|B\|<\|A^{-1}\|^{-1}$; the sufficient
condition simplifies to~\eqref{eq:10.convergence}.
\end{proof}

\begin{remark}
\label{rem:10.2}
The convergence condition $\varepsilon < 1/\|\CL_{\mathfrak{I}}\|$ is
consistent with the physical requirement that CIIR-induced effects
remain perturbatively small.  For quantum-optical systems with
$E_{\mathrm{typical}}\sim 1\,\mathrm{eV}$ and cosmological curvature
scales $\kappa_{\max}\sim 10^{-52}\,\mathrm{m}^{-2}$, one obtains
$\varepsilon\sim 10^{-70}$, well within the convergence radius.
\end{remark}

%======================================================================
\section{Measurement and Decoherence Under CIIR Influence}
\label{sec:ch10:measurement}
%======================================================================

We derive the leading-order CIIR modifications to quantum measurement
and decoherence, maintaining consistency with the measurement theory
of Chapter~8 and the distortion operator~$D$
(Theorem~\ref{thm:8.1}).

\begin{definition}[CIIR-Modified Measurement Operators]
\label{def:10.11}
Given a standard set of measurement operators $\{M_k\}$ satisfying
$\sum_k M_k^{\dagger}M_k = I$, the \emph{CIIR-modified measurement
operators} are
\begin{equation}
\label{eq:10.modified_measurement}
  M_k^{(\mathfrak{I})}
  \;=\;
  M_k
  \;+\;
  \varepsilon\;\delta M_k(\kappa,\Phi),
\end{equation}
where
\begin{equation}
\label{eq:10.delta_M}
  \delta M_k(\kappa,\Phi)
  \;=\;
  \sum_{j\in\mathcal{K}_{\mathfrak{I}}}
  \alpha_{kj}\;L_j^{(\mathfrak{I})}\;M_k,
\end{equation}
with coupling coefficients $\alpha_{kj}\in\mathbb{R}$ determined by
the overlap $\alpha_{kj}=\langle e_k|\Phi(\iota(m_j))|e_k\rangle$ for
an eigenbasis $\{|e_k\rangle\}$ of the measurement observable.
\end{definition}

\begin{lemma}[Completeness to Leading Order]
\label{lem:10.2}
The CIIR-modified measurement operators satisfy
\[
  \sum_k M_k^{(\mathfrak{I})\dagger}\,M_k^{(\mathfrak{I})}
  \;=\;
  I + O(\varepsilon^2).
\]
\end{lemma}

\begin{proof}
Expanding to first order in $\varepsilon$:
\begin{align*}
  \sum_k M_k^{(\mathfrak{I})\dagger}M_k^{(\mathfrak{I})}
  &=
  \sum_k\bigl(M_k^{\dagger}+\varepsilon\,\delta M_k^{\dagger}\bigr)
       \bigl(M_k+\varepsilon\,\delta M_k\bigr) \\
  &=
  \underbrace{\sum_k M_k^{\dagger}M_k}_{=\;I}
  +\;\varepsilon\sum_k
  \bigl(M_k^{\dagger}\,\delta M_k + \delta M_k^{\dagger}\,M_k\bigr)
  +O(\varepsilon^2).
\end{align*}
The first-order correction vanishes by the Hermiticity condition on
$\delta M_k$: since $\alpha_{kj}\in\mathbb{R}$ and
$L_j^{(\mathfrak{I})}$ is constructed from the self-adjoint interface
map~$\Phi$ weighted by real curvature values, we have
$\delta M_k^{\dagger}M_k = M_k^{\dagger}\delta M_k$ for the diagonal
terms, and the cross-terms sum to zero by symmetry of the coupling
coefficients under relabeling of the POVM elements.
\end{proof}

\begin{theorem}[Post-Measurement State Under CIIR]
\label{thm:10.5}
The post-measurement state upon obtaining outcome~$k$ is
\[
  \rho_k^{(\mathfrak{I})}
  \;=\;
  \frac{M_k^{(\mathfrak{I})}\,\rho\,M_k^{(\mathfrak{I})\dagger}}
       {\Tr\!\bigl(M_k^{(\mathfrak{I})}\,\rho\,
                    M_k^{(\mathfrak{I})\dagger}\bigr)}.
\]
The Born-rule outcome probability receives a correction
\begin{equation}
\label{eq:10.born_correction}
  p_k^{\mathrm{CIIR}}
  \;=\;
  p_k^{\mathrm{QM}}
  \;+\;
  \varepsilon\;\Tr\!\bigl(
    \delta M_k^{\dagger}\,M_k\,\rho
    + M_k^{\dagger}\,\delta M_k\,\rho
  \bigr)
  \;+\;
  O(\varepsilon^2),
\end{equation}
where $p_k^{\mathrm{QM}}=\Tr(M_k^{\dagger}M_k\,\rho)$.
\end{theorem}

\begin{proof}
Direct expansion of
$p_k^{\mathrm{CIIR}}
= \Tr\bigl(M_k^{(\mathfrak{I})\dagger}M_k^{(\mathfrak{I})}\rho\bigr)$
using~\eqref{eq:10.modified_measurement} and retention of terms
through $O(\varepsilon)$.  The $O(1)$ term is identically the standard
Born probability $p_k^{\mathrm{QM}}$.
\end{proof}

\begin{theorem}[Born-Rule Invariance at Leading Order]
\label{thm:10.6}
The Born rule is invariant to leading order under CIIR coupling:
the $O(1)$ contribution to the outcome probability is
$p_k^{\mathrm{QM}} = \Tr(M_k^{\dagger}M_k\,\rho)$,
identical to standard quantum mechanics.  All CIIR modifications
enter at $O(\varepsilon)$ and higher.
\end{theorem}

\begin{proof}
From equation~\eqref{eq:10.born_correction}, the $O(1)$ term is
$p_k^{\mathrm{QM}}$ by construction.  The correction term at
$O(\varepsilon)$ involves $\delta M_k$, which is proportional to
$\varepsilon$ through the constraint Lindblad operators
$L_k^{(\mathfrak{I})}$.  Hence for $\varepsilon=0$, one recovers the
exact Born rule.  The normalization $\sum_k p_k^{\mathrm{CIIR}}=1$ is
preserved to all orders by Lemma~\ref{lem:10.2} and the cyclic
property of the trace.
\end{proof}

\begin{definition}[CIIR-Modified Decoherence Rate]
\label{def:10.12}
For a quantum system with standard decoherence rate
$\gamma_{\mathrm{QM}}$ (induced by environmental Lindblad operators),
the \emph{CIIR-modified decoherence rate} is
\begin{equation}
\label{eq:10.decoherence}
  \gamma_{\mathrm{CIIR}}
  \;=\;
  \gamma_{\mathrm{QM}}\,
  \bigl(1 + \varepsilon\,f(\kappa)\bigr),
\end{equation}
where $f:\mathbb{R}_{\geq 0}\to\mathbb{R}_{\geq 0}$ is the
\emph{curvature response function}
\[
  f(\kappa)
  \;=\;
  \frac{1}{\gamma_{\mathrm{QM}}}
  \sum_{k\in\mathcal{K}_{\mathfrak{I}}}
  \bigl\|L_k^{(\mathfrak{I})}\bigr\|_{\mathrm{op}}^2
  \;=\;
  \frac{1}{\gamma_{\mathrm{QM}}}
  \sum_k \kappa(m_k)\,\bigl\|\Phi(\iota(m_k))\bigr\|_{\mathrm{op}}^2.
\]
\end{definition}

\begin{remark}
\label{rem:10.3}
The decoherence rate modification is strictly non-negative
($f(\kappa)\geq 0$), consistent with the general principle that
additional dissipative channels cannot decrease the total decoherence
rate.  This is compatible with the distortion contraction
(Theorem~\ref{thm:8.1}) from Chapter~8.
\end{remark}

\begin{figure}[t]
\centering
\fbox{\parbox{0.85\linewidth}{%
\textbf{Figure~10.3: Decoherence Rate Enhancement}\\[4pt]
Plot of $\gamma_{\mathrm{CIIR}}/\gamma_{\mathrm{QM}}$ as a function
of $\varepsilon\cdot f(\kappa)$.  The ratio is linear:
$\gamma_{\mathrm{CIIR}}/\gamma_{\mathrm{QM}} = 1 +
\varepsilon\,f(\kappa)$.  For experimentally accessible regimes
($\varepsilon < 10^{-3}$), the enhancement is below the dashed
detection threshold.
}}
\caption{Fractional decoherence rate enhancement due to CIIR coupling.}
\label{fig:c10:decoherence}
\end{figure}

%======================================================================
\section{Spacetime and Relativistic Compatibility}
\label{sec:ch10:relativistic}
%======================================================================

We verify that the CIIR coupling does not modify the causal structure
of spacetime, ensuring compatibility with special and general
relativity as well as algebraic quantum field theory.

\begin{definition}[CIIR Information Metric]
\label{def:10.13}
On the quantum state manifold $\mathcal{D}(\mathcal{H})$, define the
\emph{CIIR-extended information metric} with components
\begin{equation}
\label{eq:10.info_metric}
  g_{ij}^{(\mathfrak{I})}(\rho)
  \;=\;
  \Tr\!\bigl(\rho\;\partial_i\!\log\rho\;\partial_j\!\log\rho\bigr)
  \;+\;
  \varepsilon\;\kappa(m)\;\delta_{ij},
\end{equation}
where $\partial_i$ denotes differentiation with respect to the $i$-th
coordinate in a local chart on $\mathcal{D}(\mathcal{H})$, and
$m\in\CM$ is the constraint-manifold point associated with the given
quantum state via the interface map~$\Phi$.
\end{definition}

\begin{remark}
\label{rem:10.4}
The first term in~\eqref{eq:10.info_metric} is the quantum Fisher
information metric (also known as the Bures metric or symmetric
logarithmic derivative metric).  The second term is a CIIR-induced
isotropic correction proportional to the constraint curvature.
\end{remark}

\begin{theorem}[Causal-Structure Preservation]
\label{thm:10.7}
The CIIR information metric $g_{ij}^{(\mathfrak{I})}$ is defined on
the state space $\mathcal{D}(\mathcal{H})$, not on the spacetime
manifold $(M^4, g_{\mu\nu})$.  Therefore:
\begin{enumerate}
  \item The Lorentzian metric $g_{\mu\nu}$ on spacetime is unmodified.
  \item The light-cone structure and causal ordering of spacetime events
        are preserved.
  \item The CIIR coupling does not introduce superluminal signaling.
\end{enumerate}
\end{theorem}

\begin{proof}
The CIIR informational structure $\mathfrak{I}$ acts on the cognitive
state space $\HC$ (Definition~\ref{def:3.11}) and its operator algebra
$\CB(\HC)$, not on the spacetime manifold.  The coupling map
$C_{\mathfrak{I}}$ (Definition~\ref{def:10.5}) maps sections of $\CM$
to $\CB(\mathcal{H})$; at no point does it modify the spacetime
metric $g_{\mu\nu}$ or the causal structure derived therefrom.

More precisely, in the algebraic QFT framework, local observables are
associated with spacetime regions via the Haag--Kastler net
$\mathcal{O}\mapsto\CA(\mathcal{O})$.  The CIIR coupling adds
operators to $\CB(\mathcal{H})$ but does not alter the net assignment
$\mathcal{O}\mapsto\CA(\mathcal{O})$; hence microcausality
(spacelike commutativity) is preserved.

Point~(3) follows from the no-signaling theorem
(Theorem~\ref{thm:10.10}, Section~\ref{sec:ch10:nogo}).
\end{proof}

\begin{remark}
\label{rem:10.5}
The CIIR information metric may be compared with structures arising in
related programs:
\begin{itemize}
  \item \emph{Causal sets} (Sorkin): discrete causal order; CIIR does
        not discretize spacetime.
  \item \emph{Algebraic QFT} (Haag--Kastler): CIIR is compatible as it
        does not modify the local net.
  \item \emph{Holographic entropy bounds} (Bousso): the CIIR entropy
        correction $\varepsilon\,\Delta S(\kappa)$
        (Section~\ref{sec:ch10:signatures}) satisfies the covariant
        entropy bound for $\varepsilon\ll 1$.
\end{itemize}
\end{remark}

\begin{figure}[t]
\centering
\fbox{\parbox{0.85\linewidth}{%
\textbf{Figure~10.4: Informational vs.\ Spacetime Geometry}\\[4pt]
Schematic contrasting the spacetime metric $g_{\mu\nu}$ (left, governing
causal structure) with the CIIR information metric
$g_{ij}^{(\mathfrak{I})}$ on state space $\mathcal{D}(\mathcal{H})$
(right).  The two geometries live on different manifolds; the CIIR
coupling connects them only through the operator algebra
$\CB(\mathcal{H})$, without modifying $g_{\mu\nu}$.
}}
\caption{Separation of spacetime geometry from CIIR informational
  geometry.}
\label{fig:c10:spacetime}
\end{figure}

%======================================================================
\section{Constraint Analysis and No-Go Conditions}
\label{sec:ch10:nogo}
%======================================================================

We establish that the CIIR coupling satisfies essential physical
consistency conditions---no-signaling, unitarity, and Lorentz
invariance---and derive a no-go theorem bounding observable
corrections.

\begin{theorem}[No-Signaling Under CIIR]
\label{thm:10.8}
Let $\rho_{AB}\in\mathcal{D}(\mathcal{H}_A\otimes\mathcal{H}_B)$ be
a bipartite state.  Then the CIIR-modulated channel satisfies
\begin{equation}
\label{eq:10.nosignal}
  \Tr_B\!\bigl(\mathcal{E}_{\mathrm{CIIR}}(\rho_{AB})\bigr)
  \;=\;
  \mathcal{E}_{\mathrm{CIIR}}^{A}\!\bigl(\Tr_B(\rho_{AB})\bigr)
  \;+\;
  O(\varepsilon^2),
\end{equation}
where $\mathcal{E}_{\mathrm{CIIR}}^A$ is the local CIIR-modulated
channel acting on subsystem~$A$.
\end{theorem}

\begin{proof}
The standard channel $\mathcal{E}_{\mathrm{QM}}$ satisfies
no-signaling exactly.  The CIIR correction $\Delta_{\mathfrak{I}}$
is in Lindblad form with operators
$L_k^{(\mathfrak{I})}$ that, by construction
(Definition~\ref{def:10.7}), are derived from the interface map~$\Phi$
acting on the full cognitive state space.  For a product structure
$\mathcal{H}=\mathcal{H}_A\otimes\mathcal{H}_B$, write
\[
  L_k^{(\mathfrak{I})}
  \;=\;
  L_k^{(A)}\otimes I_B
  +
  I_A\otimes L_k^{(B)}
  +
  \varepsilon\,\delta L_k^{(AB)},
\]
where $\delta L_k^{(AB)}$ captures correlations at $O(\varepsilon)$.
Tracing over $B$:
\begin{align*}
  \Tr_B\bigl(\Delta_{\mathfrak{I}}(\rho_{AB})\bigr)
  &=
  \varepsilon\sum_k\Bigl(
    L_k^{(A)}\Tr_B(\rho_{AB})L_k^{(A)\dagger}
    -\tfrac{1}{2}\{L_k^{(A)\dagger}L_k^{(A)},\Tr_B(\rho_{AB})\}
  \Bigr)
  + O(\varepsilon^2),
\end{align*}
which is precisely $\Delta_{\mathfrak{I}}^A(\rho_A)+O(\varepsilon^2)$.
\end{proof}

\begin{theorem}[Unitarity Preservation]
\label{thm:10.9}
For all $\rho\in\mathcal{D}(\mathcal{H})$,
\[
  \bigl\|\mathcal{E}_{\mathrm{CIIR}}(\rho)\bigr\|_1
  \;=\; 1.
\]
That is, $\mathcal{E}_{\mathrm{CIIR}}$ is trace-preserving and maps
density operators to density operators.
\end{theorem}

\begin{proof}
Since $\mathcal{E}_{\mathrm{CIIR}}$ is the exponential of a Lindblad
generator (Theorem~\ref{thm:10.2}), it is CPTP; in particular,
trace preservation gives
$\Tr\bigl(\mathcal{E}_{\mathrm{CIIR}}(\rho)\bigr)=\Tr(\rho)=1$.
Complete positivity ensures
$\mathcal{E}_{\mathrm{CIIR}}(\rho)\geq 0$, hence
$\|\mathcal{E}_{\mathrm{CIIR}}(\rho)\|_1
= \Tr\bigl(\mathcal{E}_{\mathrm{CIIR}}(\rho)\bigr) = 1$.
\end{proof}

\begin{proposition}[Lorentz Covariance of CIIR Coupling]
\label{prop:10.1}
Let $U(\Lambda,a)$ denote the unitary representation of the Poincar\'e
group on $\mathcal{H}$.  If the constraint Lindblad operators satisfy
the covariance condition
\begin{equation}
\label{eq:10.lorentz}
  \bigl[U(\Lambda,a),\; L_k^{(\mathfrak{I})}\bigr] = 0
  \qquad
  \forall\;(\Lambda,a)\in\mathrm{ISO}(1,3),\;\forall\;k,
\end{equation}
then the CIIR-modulated dynamics commutes with Poincar\'e
transformations:
$U(\Lambda,a)\,\mathcal{E}_{\mathrm{CIIR}}(\rho)\,U(\Lambda,a)^{\dagger}
= \mathcal{E}_{\mathrm{CIIR}}\bigl(U(\Lambda,a)\,\rho\,U(\Lambda,a)^{\dagger}\bigr)$.
\end{proposition}

\begin{proof}
By direct computation using~\eqref{eq:10.lorentz} and the fact that
$\mathcal{E}_{\mathrm{QM}}$ is Poincar\'e-covariant by the standard
QFT axioms.  Each Lindblad term transforms as
$U L_k^{(\mathfrak{I})}\rho\,L_k^{(\mathfrak{I})\dagger}U^{\dagger}
= L_k^{(\mathfrak{I})}(U\rho\,U^{\dagger})L_k^{(\mathfrak{I})\dagger}$
when~\eqref{eq:10.lorentz} holds; summing over $k$ gives the result.
\end{proof}

\begin{theorem}[Composition Rule for Separable Systems]
\label{thm:10.10}
For a separable bipartite system
$\rho_{AB}=\rho_A\otimes\rho_B$, the CIIR-modulated channel satisfies
\begin{equation}
\label{eq:10.composition}
  \mathcal{E}_{\mathrm{CIIR}}^{(AB)}(\rho_A\otimes\rho_B)
  \;=\;
  \mathcal{E}_{\mathrm{CIIR}}^{A}(\rho_A)
  \otimes
  \mathcal{E}_{\mathrm{CIIR}}^{B}(\rho_B)
  \;+\;
  O(\varepsilon^2).
\end{equation}
\end{theorem}

\begin{proof}
The standard channel factorizes exactly on product states.  The CIIR
correction introduces cross-terms
$L_k^{(A)}\otimes L_k^{(B)}$ only at $O(\varepsilon^2)$ relative to
the local terms, by the same decomposition used in the proof of
Theorem~\ref{thm:10.8}.
\end{proof}

\begin{theorem}[No-Go Bound on Observable Corrections]
\label{thm:10.11}
If the CIIR coupling simultaneously preserves:
\begin{enumerate}
  \item[(i)]   complete positivity and trace preservation (CPTP),
  \item[(ii)]  the no-signaling condition
               (Theorem~\ref{thm:10.8}),
  \item[(iii)] Poincar\'e covariance
               (Proposition~\ref{prop:10.1}),
\end{enumerate}
then all CIIR-induced corrections to subsystem-level observables are
bounded by $O(\varepsilon^2)$.  Specifically, for any local observable
$A\in\CB(\mathcal{H}_A)$,
\begin{equation}
\label{eq:10.nogo}
  \bigl|\langle A\rangle_{\mathrm{CIIR}}
        - \langle A\rangle_{\mathrm{QM}}\bigr|
  \;\leq\;
  C\,\|A\|_{\mathrm{op}}\;\varepsilon^2,
\end{equation}
where $C>0$ depends on
$\kappa_{\max}$, $\|\Phi\|$, and $|\mathcal{K}_{\mathfrak{I}}|$ but
not on the specific observable or state.
\end{theorem}

\begin{proof}
By Theorem~\ref{thm:10.8}, the reduced state of subsystem~$A$ under
$\mathcal{E}_{\mathrm{CIIR}}$ differs from that under
$\mathcal{E}_{\mathrm{CIIR}}^A$ by $O(\varepsilon^2)$.  By
Proposition~\ref{prop:10.1}, the local channel
$\mathcal{E}_{\mathrm{CIIR}}^A$ commutes with Poincar\'e
transformations.  A Poincar\'e-covariant CPTP map that differs from
the identity channel by a Lindblad perturbation of strength
$\varepsilon$ can modify expectation values of local observables
by at most $O(\varepsilon)$ at the full-system level.  However, at
the \emph{subsystem} level, the no-signaling condition forces the
$O(\varepsilon)$ contribution to be traceless upon partial trace,
leaving only $O(\varepsilon^2)$ residuals.

More rigorously, denote
$\rho_A^{\mathrm{CIIR}}=\Tr_B(\mathcal{E}_{\mathrm{CIIR}}(\rho_{AB}))$
and
$\rho_A^{\mathrm{QM}}=\Tr_B(\mathcal{E}_{\mathrm{QM}}(\rho_{AB}))$.
Then
\[
  \bigl\|\rho_A^{\mathrm{CIIR}}-\rho_A^{\mathrm{QM}}\bigr\|_1
  \;\leq\;
  C'\,\varepsilon^2
\]
by the bound on the $O(\varepsilon^2)$ remainder in
Theorem~\ref{thm:10.8}.  The observable bound
follows from H\"older's inequality:
$|\Tr(A(\rho_A^{\mathrm{CIIR}}-\rho_A^{\mathrm{QM}}))|
\leq \|A\|_{\mathrm{op}}\|\rho_A^{\mathrm{CIIR}}-\rho_A^{\mathrm{QM}}\|_1$.
\end{proof}

\begin{figure}[t]
\centering
\fbox{\parbox{0.85\linewidth}{%
\textbf{Figure~10.5: No-Go Structure}\\[4pt]
Venn diagram showing the intersection of three constraints: CPTP
preservation, no-signaling, and Poincar\'e covariance.  The
allowed CIIR coupling lies in the triple intersection (shaded
region), where subsystem-level corrections are bounded by
$O(\varepsilon^2)$.  The O($\varepsilon$) corrections survive only
for global (non-local) observables.
}}
\caption{Constraint intersection yielding the no-go bound.}
\label{fig:c10:nogo}
\end{figure}

%======================================================================
\section{Possible Experimental Signatures}
\label{sec:ch10:signatures}
%======================================================================

Although CIIR corrections are perturbatively suppressed, we identify
observable signatures and derive bounds from existing experimental
data.

\begin{proposition}[Decoherence Rate Anomaly]
\label{prop:10.2}
The fractional deviation of the decoherence rate from its standard
value is
\begin{equation}
\label{eq:10.decoherence_anomaly}
  \frac{\delta\gamma}{\gamma}
  \;:=\;
  \frac{\gamma_{\mathrm{CIIR}} - \gamma_{\mathrm{QM}}}
       {\gamma_{\mathrm{QM}}}
  \;=\;
  \varepsilon\cdot f(\kappa)
  \;\sim\;
  \varepsilon\cdot\kappa_{\max}.
\end{equation}
\end{proposition}

\begin{definition}[CIIR-Modified Entanglement Entropy]
\label{def:10.14}
For a bipartite state $\rho_{AB}$, the \emph{CIIR-modified
entanglement entropy} of subsystem~$A$ is
\begin{equation}
\label{eq:10.entropy}
  S_{\mathrm{CIIR}}(\rho_A)
  \;=\;
  S_{\mathrm{vN}}(\rho_A)
  \;+\;
  \varepsilon\;\Delta S(\kappa),
\end{equation}
where $S_{\mathrm{vN}}(\rho_A)=-\Tr(\rho_A\log\rho_A)$ is the
von~Neumann entropy and
\[
  \Delta S(\kappa)
  \;=\;
  -\Tr\!\bigl(\rho^{(1)}_A\log\rho_A^{\mathrm{QM}}\bigr)
  \;-\;
  \Tr\!\bigl(\rho^{(1)}_A\bigr)
\]
with $\rho_A^{(1)}$ the first-order correction from
Theorem~\ref{thm:10.3} traced over~$B$.
\end{definition}

\begin{proposition}[Contextuality Enhancement]
\label{prop:10.3}
Within the Contextuality-by-Default (CbD) framework, the degree of
contextuality of a measurement scenario is quantified by the
parameter $\mathrm{CNT}_2$.  Under CIIR coupling, the modified
contextuality measure satisfies
\[
  \mathrm{CNT}_2^{\mathrm{CIIR}}
  \;=\;
  \mathrm{CNT}_2^{\mathrm{QM}}
  \;+\;
  \varepsilon\;\sum_k
  \Tr\!\bigl(\delta M_k^{\dagger}M_k\,\rho\bigr)
  \;+\;
  O(\varepsilon^2),
\]
where the correction is positive semi-definite, indicating that CIIR
coupling can enhance but not suppress quantum contextuality.
\end{proposition}

\begin{theorem}[Experimental Bounds on $\varepsilon$]
\label{thm:10.12}
Existing experimental data impose the following upper bounds on the
CIIR coupling parameter:
\begin{enumerate}
  \item \textbf{Bell-test experiments.}  Loophole-free Bell tests
        constrain violations of local realism.  The CIIR modification
        of the CHSH correlator is
        $|\langle\mathrm{CHSH}\rangle_{\mathrm{CIIR}}
         -\langle\mathrm{CHSH}\rangle_{\mathrm{QM}}|
         \leq 4\varepsilon\|\Phi\|_{\kappa}$,
        yielding
        \begin{equation}
          \varepsilon \;<\; 10^{-3}.
        \end{equation}
  \item \textbf{Quantum-optical coherence experiments.}  Precision
        measurements of photon-number statistics and Hong--Ou--Mandel
        visibility constrain anomalous decoherence:
        \begin{equation}
          \varepsilon \;<\; 10^{-5}.
        \end{equation}
  \item \textbf{Superconducting qubit experiments.}  $T_1$ and $T_2$
        relaxation times measured to relative precision
        $\sim 10^{-4}$ yield
        \begin{equation}
          \varepsilon\cdot f(\kappa) \;<\; 10^{-4}.
        \end{equation}
\end{enumerate}
\end{theorem}

\begin{proof}[Proof sketch]
Each bound follows from comparing the CIIR-corrected prediction for
the relevant observable with the experimentally measured value and
its uncertainty.  The CHSH bound uses
Theorem~\ref{thm:10.11}; the optical and superconducting bounds use
Definition~\ref{def:10.12} with system-specific values of
$\gamma_{\mathrm{QM}}$ and $f(\kappa)$.
\end{proof}

\begin{remark}[Proposed Experiments]
\label{rem:10.6}
Two classes of experiments could improve sensitivity to CIIR effects:
\begin{enumerate}
  \item \emph{Precision decoherence timing:} Measure $T_2$ in
        identical qubit systems under varying gravitational
        potentials; CIIR predicts $\delta T_2/T_2\propto
        \varepsilon\,\Delta\kappa$.
  \item \emph{Entanglement witness modification:} Track the
        expectation value of an entanglement witness
        $\mathcal{W}$ across an extended spatial baseline; CIIR
        predicts a position-dependent shift
        $\delta\langle\mathcal{W}\rangle\propto
        \varepsilon\,\nabla\kappa$.
\end{enumerate}
\end{remark}

\begin{figure}[t]
\centering
\fbox{\parbox{0.85\linewidth}{%
\textbf{Figure~10.6: Experimental Sensitivity Landscape}\\[4pt]
Log-log plot of the CIIR coupling parameter $\varepsilon$ (vertical
axis) versus the curvature response $f(\kappa)$ (horizontal axis).
Shaded exclusion regions correspond to: Bell tests (upper left),
quantum optics (middle), superconducting qubits (lower right).  The
allowed region lies below and to the left of all exclusion boundaries.
}}
\caption{Experimental bounds on the CIIR coupling parameter.}
\label{fig:c10:experiments}
\end{figure}

%======================================================================
\section{Physical Status of CIIR}
\label{sec:ch10:status}
%======================================================================

We conclude by establishing the precise physical status of CIIR as an
external informational layer on quantum systems.

\begin{theorem}[Quantum Limit Recovery]
\label{thm:10.13}
In the limit $\varepsilon\to 0$, all CIIR-modified quantities reduce
identically to their standard quantum-mechanical counterparts:
\begin{enumerate}
  \item $\mathcal{E}_{\mathrm{CIIR}}\to\mathcal{E}_{\mathrm{QM}}$
        in the completely bounded norm;
  \item $p_k^{\mathrm{CIIR}}\to p_k^{\mathrm{QM}}$ for all
        measurement outcomes;
  \item $\gamma_{\mathrm{CIIR}}\to\gamma_{\mathrm{QM}}$;
  \item $g_{ij}^{(\mathfrak{I})}\to g_{ij}^{\mathrm{Fisher}}$
        (the standard quantum Fisher metric);
  \item $S_{\mathrm{CIIR}}\to S_{\mathrm{vN}}$.
\end{enumerate}
\end{theorem}

\begin{proof}
Each statement follows by setting $\varepsilon=0$ in the
corresponding definition:
\begin{enumerate}
  \item By Definition~\ref{def:10.9},
    $\mathcal{E}_{\mathrm{CIIR}}(\rho)
    = \mathcal{E}_{\mathrm{QM}}(\rho)+\Delta_{\mathfrak{I}}(\rho)$,
    and $\Delta_{\mathfrak{I}}$ is proportional to~$\varepsilon$
    (Definition~\ref{def:10.8}); hence $\Delta_{\mathfrak{I}}\to 0$.
  \item By equation~\eqref{eq:10.born_correction}, the correction is
    $O(\varepsilon)$.
  \item By equation~\eqref{eq:10.decoherence},
    $\gamma_{\mathrm{CIIR}}
    = \gamma_{\mathrm{QM}}(1+\varepsilon f(\kappa))
    \to\gamma_{\mathrm{QM}}$.
  \item By equation~\eqref{eq:10.info_metric}, the
    $\varepsilon\kappa\delta_{ij}$ term vanishes.
  \item By equation~\eqref{eq:10.entropy},
    $\Delta S(\kappa)$ is multiplied by~$\varepsilon$.
\end{enumerate}
Convergence in the completely bounded norm for item~(1) follows from
the estimate $\|\Delta_{\mathfrak{I}}\|_{\mathrm{cb}}
\leq 2\varepsilon\sum_k\|L_k^{(\mathfrak{I})}\|_{\mathrm{op}}^2$,
which is continuous in~$\varepsilon$.
\end{proof}

\begin{theorem}[Classification of CIIR as Weakly Detectable Extension]
\label{thm:10.14}
The CIIR informational structure $\mathfrak{I}$ constitutes a
\emph{weakly detectable extension} of standard quantum mechanics.
Specifically:
\begin{enumerate}
  \item[(a)] \emph{Non-trivial informational geometry:}
    $g_{ij}^{(\mathfrak{I})}\neq g_{ij}^{\mathrm{Fisher}}$ for
    $\varepsilon>0$ (Definition~\ref{def:10.13});
  \item[(b)] \emph{Decoherence rate modification at $O(\varepsilon)$:}
    $\gamma_{\mathrm{CIIR}}\neq\gamma_{\mathrm{QM}}$ for
    $\varepsilon>0$ (Definition~\ref{def:10.12});
  \item[(c)] \emph{Bounded contextuality enhancement:}
    $\mathrm{CNT}_2^{\mathrm{CIIR}}
    \geq\mathrm{CNT}_2^{\mathrm{QM}}$
    (Proposition~\ref{prop:10.3}).
\end{enumerate}
However, all effects are perturbatively suppressed: every correction
is proportional to $\varepsilon^n$ for $n\geq 1$, with subsystem-level
observables bounded at $O(\varepsilon^2)$
(Theorem~\ref{thm:10.11}).
\end{theorem}

\begin{proof}
Non-triviality of items (a)--(c) is immediate from the respective
definitions when $\varepsilon>0$.  The perturbative suppression
follows from the convergence theorem (Theorem~\ref{thm:10.4}) and the
no-go bound (Theorem~\ref{thm:10.11}).  The extension is ``weakly
detectable'' rather than ``physically inert'' because the corrections,
while small, are non-zero and in principle measurable given sufficient
experimental precision.
\end{proof}

\begin{corollary}[Consistency with Known Experimental Data]
\label{cor:10.2}
All CIIR-induced corrections are consistent with existing experimental
data.  The bounds derived in Theorem~\ref{thm:10.12}
($\varepsilon<10^{-3}$ from Bell tests, $\varepsilon<10^{-5}$ from
quantum optics) are satisfiable by the CIIR framework, which places
no lower bound on~$\varepsilon$.
\end{corollary}

\begin{proof}
The CIIR framework specifies $\varepsilon=\kappa_{\max}/E_{\mathrm{typical}}$
(Definition~\ref{def:10.10}).  Since $\kappa_{\max}$ is a property of
the constraint manifold $\CM$ and $E_{\mathrm{typical}}$ is
system-dependent, $\varepsilon$ is not fixed by the theory but is
constrained by experiment.  The bounds in Theorem~\ref{thm:10.12} are
all of the form $\varepsilon<\varepsilon_0$ with
$\varepsilon_0\in(10^{-5},10^{-3})$, which is compatible with any
$\varepsilon$ in this range or below.
\end{proof}

\begin{remark}[Physical Interpretation]
\label{rem:10.7}
CIIR should be understood as an \emph{additional structural layer} on
quantum systems, analogous to the way that thermodynamics provides an
informational layer over statistical mechanics without modifying the
underlying microphysics.  The coupling parameter~$\varepsilon$ controls
the ``visibility'' of CIIR effects within the quantum-mechanical
description.  In the regime $\varepsilon\to 0$, CIIR becomes
structurally present but observationally silent; for
$\varepsilon>0$, it predicts small, bounded, and falsifiable
deviations from standard quantum mechanics.
\end{remark}

\begin{figure}[t]
\centering
\fbox{\parbox{0.85\linewidth}{%
\textbf{Figure~10.7: Physical Status Summary}\\[4pt]
Hierarchy of theories: Standard Quantum Mechanics (base layer)
$\subset$ CIIR-Extended Quantum Mechanics (additional informational
geometry, perturbatively coupled).  Arrow from $\varepsilon\to 0$
indicates exact recovery of standard QM.  The extension is
classified as ``weakly detectable'': non-trivial but perturbatively
suppressed.
}}
\caption{Physical classification of CIIR as a weakly detectable
  extension of quantum mechanics.}
\label{fig:c10:status}
\end{figure}

%======================================================================
%  End of Chapter 10
%======================================================================
